%% kapa: exam=1 round=29 year=2018 period=2 form=A subject=accounting subject_ko=회계학 seq=4 number=44
%% source: Q-Net 감정평가사 기출문제 (cst00309) artlSeq=5202808 file=extracted/1008/5202808/2018년 제29회 감정평가사 1차시험 기출문제/2018년 제29회 감정평가사 1차시험 기출문제/2교시 A형.pdf
%% conversion: kordoc 4.24.51 -> tools/gen1.py (auto); review=proofread-sonnet; flags=none
\begin{question}{44}
다음은 (주)감평의 20x1년도 재고자산 거래와 관련된 자료이다.
\begin{framedtext}
\centering
\begin{tabular}{cccr}
\uline{일자} & \uline{적요} & \uline{수량} & \uline{단가} \\
1월 1일 & 기초재고 & 100개 & ￦90 \\
3월 9일 & 매입 & 200개 & 150 \\
5월 16일 & 매출 & 150개 & \\
8월 20일 & 매입 & 50개 & 200 \\
10월 25일 & 매입 & 50개 & 220 \\
11월 28일 & 매출 & 200개 & \\
\end{tabular}
\end{framedtext}
다음 설명 중 옳지 않은 것은? \par
\begin{choices}
  \item 실지재고조사법을 적용하여 선입선출법을 사용할 경우 기말재고자산 금액은 ￦11,000이다.
  \item 실지재고조사법을 적용하여 가중평균법을 사용할 경우 매출원가는 ￦52,500이다.
  \item 선입선출법을 사용할 경우보다 가중평균법을 사용할 때 당기순이익이 더 작다.
  \item 가중평균법을 사용할 경우, 실지재고조사법을 적용하였을 때보다 계속기록법을 적용하였을 때 당기순이익이 더 크다.
  \item 선입선출법을 사용할 경우, 계속기록법을 적용하였을 때보다 실지재고조사법을 적용하였을 때 매출원가가 더 크다.
\end{choices}
\end{question}
\answer{5}
